% Zdroj: PDF priklady-leto-2024.pdf, fyzická str. 2-3. Značka: společné okruhy. Zařazeno: I4.
\section{Semafor}
\szzotazka{léto 2024}{2}{Semafor (P/V, race condition, deadlock)}

\noindent\textbf{Zadání.}\par
Knihovna \texttt{System.Threading} jazyka C\# obsahuje třídu \texttt{Semaphore}
s~metodami odpovídajícími operacím klasického semaforu, $P$ (\texttt{WaitOne})
a~$V$ (\texttt{Release}). Objekt vytvořený pomocí \texttt{new Semaphore(0, 1)} odpovídá
binárnímu semaforu inicializovanému nulou. Pokud více než jedno vlákno čeká na tentýž
semafor, pořadí, ve kterém budou vlákna spouštěna po uvolnění semaforu, není specifikováno.

Program vpravo má produkovat výstup \texttt{OneTwoThreeFourFive}, přitom liché části
výstupu mají být vypisovány funkcí \texttt{f1} a~sudé části funkcí \texttt{f2}, přičemž
tyto funkce běží v~různých vláknech.

Synchronizace, implementovaná pomocí semaforu, téměř funguje, ale obsahuje časově závislé
chyby (race conditions).

\begin{lstlisting}[language={[Sharp]C},numbers=left]
class Program
{
    private static Semaphore sem;

    private static void f1()
    {
        Console.Write("One");
        sem.Release();
        sem.WaitOne();
        Console.Write("Three");
        sem.Release();
        sem.WaitOne();
        Console.Write("Five");
    }

    private static void f2()
    {
        sem.WaitOne();
        Console.Write("Two");
        sem.Release();
        sem.WaitOne();
        Console.Write("Four");
        sem.Release();
    }

    static void Main(string[] args)
    {
        sem = new Semaphore(0, 1);
        Thread t1 = new Thread(f1);
        Thread t2 = new Thread(f2);
        t1.Start();
        t2.Start();
        t1.Join();
        t2.Join();
    }
}
\end{lstlisting}

\begin{enumerate}
\item Popište scénář, který vede k~jinému výstupu než \texttt{OneTwoThreeFourFive}, nebo
  končí uváznutím (deadlock). Scénář zapište jako posloupnost čísel řádek, přičemž každé
  číslo reprezentuje \textit{ukončení} příkazu na dané řádce. Pouze vnitřky funkcí
  \texttt{f1} a~\texttt{f2} jsou relevantní.
\item Existovaly by v~tomto kódu časově závislé chyby i~v~případě, že by knihovna C\#
  garantovala FIFO pořadí spouštění vláken čekajících ve \texttt{WaitOne}?
\item Napište řešení, které neobsahuje časově závislé chyby a~spolehlivě vypisuje
  požadovaný výstup \texttt{OneTwoThreeFourFive}. Použijte přitom dva semafory a~pouze
  jejich funkce \texttt{Release()} a~\texttt{WaitOne()}. Řešení nesmí používat žádné jiné
  proměnné nebo objekty sdílené mezi vlákny kromě těchto dvou semaforů. Volání
  \texttt{Console.Write} musejí samozřejmě zůstat tam, kde jsou.
\end{enumerate}

\noindent(Jazyk C\# je v~této otázce použitý pouze jako generický zástupce běžných
programovacích jazyků, otázka ani řešení s~jazykem C\# jako takovým nesouvisí.)

\medskip
\noindent\textbf{Náčrt řešení.}\par
\begin{enumerate}
\item Všechny scénáře, v~lexikografickém pořadí:
\begin{enumerate}
\item \texttt{7-8-9-10-11-12-13-deadlock(18+34): OneThreeFive}
\item \texttt{7-8-9-10-11-18-19-20-12-13-deadlock(21+34): OneThreeTwoFive}
\item \texttt{7-8-9-10-11-18-19-20-21-22-23-12-13: OneThreeTwoFourFive}
\item \texttt{7-8-18-19-20-9-10-11-12-13-deadlock(21+34): OneTwoThreeFive}
\item \textit{\texttt{7-8-18-19-20-9-10-11-21-22-23-12-13: OneTwoThreeFourFive} -- The expected program behavior}
\item \texttt{7-8-18-19-20-21-22-23-9-10-11-12-13: OneTwoFourThreeFive}
\end{enumerate}
Postačující odpovědí je kterýkoli z~výše uvedených scénářů kromě (e).

\item Ano. Všechny výše uvedené scénáře fungují i~pro FIFO čekání.

\item Viz kód níže:
\begin{lstlisting}[language={[Sharp]C}]
private static void f1()
{
    Console.Write("One");
    semA.Release();
    semB.WaitOne();
    Console.Write("Three");
    semA.Release();
    semB.WaitOne();
    Console.Write("Five");
}

private static void f2()
{
    semA.WaitOne();
    Console.Write("Two");
    semB.Release();
    semA.WaitOne();
    Console.Write("Four");
    semB.Release();
}
\end{lstlisting}
\end{enumerate}

\medskip
\noindent\textit{Doplňující vysvětlení (nad rámec oficiálního náčrtu):} Kořen chyby je v~tom, že
jediný semafor nerozlišuje, \emph{komu} je signál určen: \texttt{f1} po \texttt{Release} (ř.~8)
pokračuje hned \texttt{WaitOne} (ř.~9) a~může si svůj vlastní token vzít dřív, než \texttt{f2}
vůbec dojde na ř.~18. Proto nepomůže ani FIFO pořadí probouzení: ve scénáři (a) na semaforu v~tu
chvíli nikdo nečeká, takže není koho probouzet ve frontě. Se dvěma semafory \texttt{f1} signalizuje
jen do \texttt{semA} a~čeká jen na \texttt{semB}, \texttt{f2} naopak; žádné vlákno tak nemůže
spotřebovat token, který samo vydalo, a~výpisy se střídají v~pevném pořadí.
